\paragraph{Transferring recognition into independent writing.}
Continuing the aligned writer for 24 epochs on its original pictured
targets increased exact development completions from 134 to 292 of
22,907. With the same passages and updates, sampling whole training
expressions from the frozen reader increased this to 378 (1.65\%).
The difference was 0.38 percentage points (source-group bootstrap
95\% interval, 0.21--0.54). Of 851 passages that the reader recognized
and the initial writer missed, the guided writer recovered 190. It
lost 52 of the initial writer's 134 exact answers overall. Of 16,797
passages missing every answer glyph, 268 were completed exactly. None of
38 cases requiring an expression unseen in training was recovered.

On 9,135 previously unused passages from 1,356 source groups, the initial,
continued and guided writers recovered 60, 125 and 163 expressions.
The guided-minus-continued difference was 0.42 percentage points (95\%
interval, 0.19--0.65). The local count baseline recovered 406, and the
supplied-expression reader recovered 398. The guided writer reproduced
112 of those 398. It solved both members of two of 3,237 pairs sharing
the same two adjacent glyphs but requiring different expressions.
Keeping only the eight nearest cells reduced 163 exact completions to
126 (66 lost, 29 gained). Removing those eight gave seven. Neither
intervention preserved a complete both-correct pair. The intact-minus-near
difference was 0.41 percentage points (95\% interval, 0.20--0.63).
These comparisons identify input dependence under masking. The task
supplies a gap and does not evaluate open-ended question answering
(Fig.~\ref{fig:writer-transfer}).
